\section{Experiments}

\subsection{Experimental Setup}

\subsubsection{Models}

We evaluate our method on a 12B-parameter open-weight language model (Gemma3-12B-IT). 
We compare against the base model (without fine-tuning) and GPT-4o as a strong closed-source baseline.

\subsubsection{Evaluation Settings}

We conduct experiments under controlled contextual conditions with varying information quality, including scenarios with no reliable information and scenarios with partially correct information. 
This setup enables evaluation of both hallucination suppression and response quality under uncertainty.

\subsubsection{Metrics}

\paragraph{Hallucination Metric}
We measure hallucination using a binary judgment indicating whether a response is supported by available information or contains unsupported or incorrect content. 
We report hallucination rate (\%) across different difficulty levels.

\paragraph{Response Quality Metrics}
We adopt an LLM-as-a-judge framework with eight evaluation dimensions, including relevance, coherence, consistency, fluency, completeness, actionability, evidence use, and conciseness. 
Most dimensions are rated on a 1--5 scale, while fluency is evaluated on a 1--3 scale. 
To ensure comparability, scores are aggregated by averaging across dimensions. 
Detailed definitions of these metrics are provided in Appendix~\ref{app:metrics}.

\subsubsection{Dataset}

We construct a controlled dataset consisting of two scenarios: (i) no reliable information (all\_wrong), where the model should abstain, and (ii) partially correct information (one\_correct), where the model should generate grounded responses. 
The dataset is derived from internal sources and is not publicly available due to privacy constraints.


\subsection{Main Results}

\subsubsection{Hallucination Reduction}

\begin{table}[h]
\centering
\caption{Hallucination rate (\%) across context sizes (lower is better).}
\label{tab:hallucination}
\begin{tabular}{cccc}
\toprule
K & GPT-4o & Base & Ours \\
\midrule
1 & 14.3\% & 24.3\% & \textbf{6.0\%} \\
3 & 16.3\% & 22.0\% & \textbf{8.3\%} \\
5 & 17.0\% & 23.3\% & \textbf{8.6\%} \\
7 & 23.0\% & 25.6\% & \textbf{8.6\%} \\
9 & 19.6\% & 26.0\% & \textbf{10.3\%} \\
\midrule
\textbf{Avg} & 18.04\% & 24.24\% & \textbf{8.36\%} \\
\bottomrule
\end{tabular}
\end{table}

As shown in Table~\ref{tab:hallucination}, our method substantially reduces hallucination compared to both the base model and GPT-4o.

On average, the hallucination rate decreases to 8.36\%, corresponding to a 65.5\% relative reduction compared to the base model and a 53.6\% reduction compared to GPT-4o.

The improvement is consistent across all context sizes, indicating robustness as contextual complexity increases.

These results suggest that explicitly modeling uncertainty and behavior enables the model to make more reliable decisions, thereby reducing hallucination in a stable and generalizable manner.

\subsubsection{Generalization Performance}

\begin{table}[h]
\centering
\caption{Hallucination rate (\%) on unseen domain dataset (lower is better).}
\label{tab:generalization}
\begin{tabular}{cccc}
\toprule
K & GPT-4o & Base & Ours \\
\midrule
1 & 2.3\% & 12.7\% & \textbf{2.3\%} \\
3 & 5.3\% & 12.3\% & \textbf{5.0\%} \\
5 & 6.7\% & 14.7\% & \textbf{5.6\%} \\
7 & 9.3\% & 18.3\% & \textbf{7.6\%} \\
9 & 8.0\% & 17.3\% & 8.6\% \\
\midrule
\textbf{Avg} & 6.32\% & 15.06\% & \textbf{5.82\%} \\
\bottomrule
\end{tabular}
\end{table}

We evaluate generalization performance under domain shift by measuring hallucination rates on unseen data.

As shown in Table~\ref{tab:generalization}, the proposed method consistently reduces hallucination compared to the base model across different values of $K$.

On average, our method achieves a hallucination rate of 5.82\%, outperforming both the base model (15.06\%) and GPT-4o (6.32\%).

Notably, performance remains competitive with GPT-4o across most settings, indicating that the learned behavior generalizes effectively beyond the training domain.

These results suggest that the proposed framework improves robustness under distribution shift while maintaining strong generalization capability.

\subsubsection{Response Quality}

\begin{table}[h]
\centering
\caption{Response quality evaluation (higher is better).}
\label{tab:response_quality}
\begin{tabular}{lcc}
\toprule
Metric & Base & Ours \\
\midrule
Relevance & 4.10 & \textbf{4.62} \\
Coherence & 3.62 & \textbf{4.42} \\
Consistency & 4.69 & \textbf{4.83} \\
Fluency & 2.86 & \textbf{2.87} \\
Completeness & 3.29 & \textbf{3.95} \\
Actionability & 2.70 & \textbf{3.32} \\
Evidence Use & 2.09 & \textbf{4.79} \\
Conciseness & 3.97 & \textbf{4.37} \\
\midrule
\textbf{Average} & 3.42 & \textbf{4.15} \\
\bottomrule
\end{tabular}
\end{table}

Response quality is evaluated across multiple dimensions following LLM-based evaluation approaches such as G-Eval~\cite{geval}.

As shown in Table~\ref{tab:response_quality}, our method consistently improves response quality across all dimensions.

The largest gain is observed in evidence use (2.09 → 4.79), indicating significantly improved grounding and transparency in generated responses.

Substantial improvements are also observed in completeness and actionability, suggesting that the model produces more informative and practically useful outputs.

Overall, the average score increases from 3.42 to 4.15, demonstrating that the proposed framework enhances response quality in a comprehensive manner.

\subsubsection{General Knowledge Preservation}

To ensure that hallucination reduction is not achieved at the expense of general capabilities,
we evaluate the model on a diverse set of benchmark datasets covering:

\begin{itemize}
    \item English knowledge (MMLU, MMLU-Pro)
    \item Chinese knowledge (CMMLU, CEVAL)
    \item Instruction following (IFEval)
    \item Mathematical reasoning (GSM8K, Math-500)
\end{itemize}

\begin{table}[h]
\centering
\caption{General knowledge evaluation (accuracy, higher is better).}
\label{tab:general_knowledge}
\begin{tabular}{l l c c}
\toprule
Domain & Dataset & Base & Ours \\
\midrule
English & MMLU     & 0.7714 & 0.7686 \\
English & MMLU-Pro & 0.6034 & 0.6000 \\
Instruction & IFEval & 0.8529 & 0.8571 \\
Chinese & CMMLU    & 0.6344 & 0.6200 \\
Chinese & CEVAL    & 0.6636 & 0.6759 \\
Math & GSM8K       & 0.9574 & 0.8842 \\
Math & Math-500    & 0.9000 & 0.8000 \\
\midrule
\textbf{Average} &  & 0.7690 & 0.7436 \\
\textbf{Average (w/o Math)} &  & 0.7051 & 0.7043 \\
\bottomrule
\end{tabular}
\\
\footnotesize{
Results are computed on a sampled subset of 3,000 instances (out of 41,512), with weighted accuracy based on dataset proportions.
}
\end{table}

The results in Table~\ref{tab:general_knowledge} show that the proposed method largely preserves general knowledge capability.

When excluding mathematical reasoning tasks, the average accuracy remains nearly unchanged (0.7051 → 0.7043), 
indicating minimal degradation in general knowledge performance.

Instruction-following ability slightly improves (IFEval: 0.8529 → 0.8571), 
suggesting that behavior-oriented training enhances controllability.

While performance decreases on mathematical reasoning tasks, 
this may be attributed to the absence of explicit optimization for multi-step reasoning in the proposed reward design.

Overall, these results demonstrate that hallucination mitigation is achieved without catastrophic forgetting \cite{EWC}, 
and the model retains its general-purpose capabilities.

\subsubsection{Summary}

Overall, the proposed method achieves a strong balance between hallucination mitigation, response quality, and general capability preservation.


\subsection{Analysis}

\subsubsection{Entropy and Overconfidence}

We analyze token-level entropy:

\[
\tilde{H} = \log_{10}(H_{\text{avg}} + \epsilon)
\]

\begin{table}[h]
\centering
\caption{Average entropy under different answerability conditions. Higher values indicate greater uncertainty.}
\label{tab:entropy}
\begin{tabular}{ccc}
\toprule
K & all\_wrong & one\_correct \\
\midrule
1 & -0.6790 & -1.5007 \\
3 & -0.8695 & -1.7862 \\
5 & -0.8462 & -1.8510 \\
7 & -0.8746 & -1.7624 \\
9 & -0.8994 & -1.7480 \\
\bottomrule
\end{tabular}
\end{table}

We analyze model uncertainty by measuring entropy under different answerability conditions.

As shown in Table~\ref{tab:entropy}, entropy is consistently higher for unanswerable cases (\textit{all\_wrong}) and lower for answerable cases (\textit{one\_correct}).

This clear separation suggests that entropy serves as an effective signal for distinguishing between answerable and unanswerable scenarios.

These results support the use of entropy as a proxy for decision boundaries, enabling the model to determine when to answer and when to abstain.

\subsubsection{Abstention Behavior}

\begin{table}[h]
\centering
\caption{Average entropy under different abstention outcomes. Higher values indicate greater uncertainty.}
\label{tab:abstention}
\begin{tabular}{ccc}
\toprule
K & Abstained & Not Abstained \\
\midrule
1 & -0.6444 & -1.1390 \\
3 & -0.7498 & -1.3853 \\
5 & -0.7753 & -1.3675 \\
7 & -0.7649 & -1.3284 \\
9 & -0.7664 & -1.5338 \\
\bottomrule
\end{tabular}
\end{table}

We analyze the relationship between entropy and abstention behavior.

As shown in Table~\ref{tab:abstention}, entropy is consistently higher for abstained cases and lower for non-abstained cases.

This suggests that correct abstention is associated with higher uncertainty, while incorrect answers tend to exhibit lower entropy, indicating overconfident predictions.

These findings support the effectiveness of the proposed entropy-based penalty in encouraging appropriate abstention behavior.